% !TeX root = ../main.tex

\chapter{系统性质与 LTI 结构}
\label{chap:systems-and-lti}

\section{系统的数学描述}
\label{sec:system-description}

系统（system）接收输入信号，按照确定的规则产生输出信号。沿用笔记中的箭头流程图，可写成
\begin{equation}
  x(t)\longrightarrow\boxed{\text{系统}}\longrightarrow y(t),
  \qquad
  x[n]\longrightarrow\boxed{\text{系统}}\longrightarrow y[n].
  \label{eq:system-flow}
\end{equation}
研究系统性质时，真正的对象是\red{框中固定的输入输出规则}，而非某一组特定的输入、输出波形。下文若写“$x$ 产生 $y$”，均指 $x$ 经过同一个待研究系统后得到 $y$。
\begin{pitfall}[不能用一组输入输出判断系统性质]
  看到某个输入恰好得到成比例的输出，并不能证明系统线性；看到某个输入平移后输出也恰好平移，也不能证明系统时不变。定义中的判据都要求对\red{任意允许的输入}成立。
\end{pitfall}

\section{线性系统}
\label{sec:linear-systems}

\begin{definition}[线性系统]
  \label{def:linear-system}
  设两组输入输出关系为
  \[
    x_1\longrightarrow\boxed{\text{系统}}\longrightarrow y_1,
    \qquad
    x_2\longrightarrow\boxed{\text{系统}}\longrightarrow y_2.
  \]
  若对任意输入 $x_1,x_2$ 和任意常数 $a,b$，都有
  \begin{equation}
    a x_1+b x_2
    \longrightarrow\boxed{\text{同一系统}}\longrightarrow
    a y_1+b y_2,
    \label{eq:linearity-test}
  \end{equation}
  则称该系统为线性系统（linear system）。这一条件也称为叠加原理（superposition principle）。
\end{definition}

\cref{eq:linearity-test}可以拆成两个缺一不可的条件：
\begin{itemize}
  \item 若 $x\longrightarrow y$，则 $a x\longrightarrow a y$，称为齐次性（homogeneity）；
  \item 若 $x_1\longrightarrow y_1$、$x_2\longrightarrow y_2$，则 $x_1+x_2\longrightarrow y_1+y_2$，称为可加性（additivity）。
\end{itemize}

\noindent\textbf{判据 }
对笔记中的常见显式系统公式，先检查：
\begin{enumerate}
  \item \red{每一项都有输入信号}，不得有与输入无关的常数项；
  \item \red{每个输入均为一次}，不得有输入的乘积、平方、绝对值、指数等非线性运算。
\end{enumerate}
第一条也可先用零输入检验。令 $a=b=0$，线性系统必满足
\begin{equation}
  0\longrightarrow\boxed{\text{线性系统}}\longrightarrow0.
  \label{eq:zero-input-zero-output}
\end{equation}
若任一快速检查不通过，就可判定系统非线性；但检查通过仍须回到\cref{eq:linearity-test} 验证齐次性和可加性。\red{证明非线性只需一个反例}，证明线性则必须对一般的 $a x_1+b x_2$ 成立。

\begin{intuition}
  先把输入拆成若干简单分量，分别求响应后再相加，必须与一次把完整输入送入系统得到相同结果。这正是以后用冲激分解研究 LTI 系统的基础。
\end{intuition}

\begin{example}[增益、微分与累加]
  系统 $y(t)=3x(t)$ 满足齐次性和可加性，因而线性。微分系统
  \[
    y(t)=\frac{\dd x(t)}{\dd t}
  \]
  也线性，因为求导可分配到线性组合上。离散差分系统
  \[
    y[n]=x[n]-x[n-1]
  \]
  与累加系统
  \[
    y[n]=\sum_{k=-\infty}^{n}x[k]
  \]
  同样满足\cref{eq:linearity-test}。这些例子说明：系统是否线性取决于\red{运算对线性组合的作用}，而不取决于系统公式看起来是否复杂。
\end{example}

\begin{example}[常数偏置与平方]
  对 $y(t)=x(t)+1$，零输入产生 $y(t)=1$，所以系统非线性。对 $y(t)=x^2(t)$，虽然零输入产生零输出，但
  \[
    2x(t)\longrightarrow4x^2(t)\ne2x^2(t),
  \]
  齐次性不成立，系统仍非线性。这也说明\red{零输入检验只能用于否定}，不能单独用于证明。
\end{example}

\begin{pitfall}[把“可加”误当成“线性”]
  只验证“输入相加时输出也相加”还不够；系统还必须允许任意幅度缩放。遇到平方、绝对值、指数或输入之间的乘积时，\red{优先检查齐次性}，通常能迅速找到矛盾。
\end{pitfall}

\section{时不变系统}
\label{sec:time-invariant-systems}

\begin{definition}[时不变系统]
  \label{def:time-invariant-system}
  若
  \[
    x(t)\longrightarrow\boxed{\text{系统}}\longrightarrow y(t),
  \]
  并且对任意实数 $t_0$ 都有
  \begin{equation}
    x(t-t_0)\longrightarrow\boxed{\text{同一系统}}\longrightarrow y(t-t_0),
    \label{eq:ct-time-invariance}
  \end{equation}
  则称该连续时间系统为时不变系统（time-invariant system）。离散时间系统的判据为：对任意整数 $n_0$，
  \begin{equation}
    x[n-n_0]\longrightarrow\boxed{\text{同一系统}}\longrightarrow y[n-n_0].
    \label{eq:dt-time-invariance}
  \end{equation}
\end{definition}

\noindent\textbf{判据 }
对 $y(t)$ 直接由若干 $x(\varphi(t))$ 构成的常见系统公式，先检查：
\begin{enumerate}
  \item \red{时间只出现在输入括号内}，公式中不应有裸露的 $t$；
  \item 除固定平移外，输入括号内应\red{保持时间尺度和方向}，即保留 $t$，不能换成 $2t$、$-t$、$t^2$ 等其他时间函数。
\end{enumerate}
这两条是笔记中的快速筛查法；含积分、求和等结构时，仍应使用定义作严格判断。时不变性的含义是：\red{输入延迟多少，输出就只延迟同样多}，除此之外形状和幅值都不改变。严格判断时比较两条路径：
\begin{enumerate}
  \item 先把输入换成 $x(t-t_0)$，再代入系统规则，得到输出 $y_1(t)$；
  \item 先由原输入得到 $y(t)$，再把输出中的\red{每一个显含时间变量}替换为 $t-t_0$，得到 $y_2(t)=y(t-t_0)$。
\end{enumerate}
若对任意输入和任意 $t_0$ 都有 $y_1=y_2$，系统才时不变。离散时间完全类似。

\begin{example}[延迟与非线性不冲突]
  对 $y(t)=x(t-1)$，输入 $x(t-t_0)$ 产生 $x(t-t_0-1)$，恰好等于 $y(t-t_0)$，所以系统时不变。系统
  \[
    y(t)=\e^{x(t+1)}
  \]
  虽然非线性且使用未来输入，但输入平移后得到 $\e^{x(t-t_0+1)}=y(t-t_0)$，因此仍然时不变。可见\red{时不变与线性、因果性彼此独立}。
\end{example}

\begin{example}[显含时间与时间缩放]
  对 $y(t)=t x(t)$，输入平移后得到 $t x(t-t_0)$，而输出平移为 $(t-t_0)x(t-t_0)$，一般不相等，所以系统时变。对 $y(t)=x(2t)$，两条路径分别给出
  \[
    x(2t-t_0),
    \qquad
    x(2t-2t_0),
  \]
  也一般不相等。因此，系统内部的时间尺度变换通常会\red{破坏时不变性}。
\end{example}

\section{因果系统}
\label{sec:causal-systems}

\begin{definition}[因果系统]
  \label{def:causal-system}
  若对任意时刻 $t_0$，只要两个输入在所有 $t\le t_0$ 上相同，相应输出在 $t_0$ 也必相同，则称连续时间系统为因果系统（causal system）。等价地，$y(t_0)$ 只能依赖输入在\red{当前与过去}的取值，不能依赖 $t>t_0$ 的输入。

  对离散时间系统，把条件改为两个输入在所有 $n\le n_0$ 上相同；此时 $y[n_0]$ 只能依赖 $x[k]$（$k\le n_0$）。
\end{definition}

\noindent\textbf{判据 }
固定输出时刻 $t_0$ 或 $n_0$，列出计算该时刻输出所需的全部输入时刻。对直接出现的 $x(\varphi(t_0))$ 或 $x[\varphi(n_0)]$，必须分别满足
\[
  \varphi(t_0)\le t_0,
  \qquad
  \varphi(n_0)\le n_0.
\]
这就是笔记所写的\red{输入括号内函数不晚于输出时刻}。积分或求和系统则检查整个读取区间；只要存在一个严格晚于输出时刻且确实影响输出的输入，就能判定系统非因果。

\begin{example}[延迟、超前与无记忆运算]
  $y(t)=x(t-1)$ 只读取过去输入，是因果系统；$y(t)=x(t+1)$ 读取未来输入，是非因果系统。$y(t)=x^2(t)$ 只依赖当前输入，因此虽非线性却因果。一般地，任何输出只依赖 $x(t)$ 或 $x[n]$ 当前值的\red{无记忆系统必然因果}。
\end{example}

\begin{example}[积分与累加的上限]
  系统
  \[
    y(t)=\int_{-\infty}^{t}x(\tau)\dd\tau,
    \qquad
    y[n]=\sum_{k=-\infty}^{n}x[k]
  \]
  只汇总当前及过去输入，因而因果。若把上限改成 $t+1$ 或 $n+1$，当前输出就会使用未来输入，系统变为非因果。对积分、求和系统，\red{先比较上限与观察时刻}通常最快。
\end{example}

\begin{example}[时间缩放系统]
  对 $y(t)=x(2t)$，在 $t_0>0$ 时需要读取 $x(2t_0)$，而 $2t_0>t_0$，所以系统非因果。只需找到这一个违反定义的时刻即可；不要求系统在每个时刻都依赖未来。
\end{example}

\begin{pitfall}[用输入的具体波形代替依赖关系]
  某个特定输入可能在未来恒为零，使非因果系统看起来没有“利用未来”。但因果性是系统映射对所有输入的性质。正确检查应关注输出公式\red{可能读取哪些输入时刻}。
\end{pitfall}

\section{无记忆系统}
\label{sec:memoryless-systems}

\begin{definition}[无记忆系统]
  \label{def:memoryless-system}
  若对任意时刻 $t_0$，输出 $y(t_0)$ 只依赖输入在同一时刻的值 $x(t_0)$，则称连续时间系统为无记忆系统（memoryless system）。离散时间系统的定义相同：$y[n_0]$ 只能依赖 $x[n_0]$，不能读取其他下标的输入样值。
\end{definition}

\noindent\textbf{判据 }
逐项比较输入与输出的时间自变量：连续时间必须只读取 $x(t)$，离散时间必须只读取 $x[n]$，也就是\red{输入输出括号内函数完全相同}。只要读取了 $x(\varphi(t))$ 且 $\varphi(t)\ne t$，或读取了 $x[\varphi(n)]$ 且 $\varphi(n)\ne n$，系统就有记忆。

这里的“记忆”并不是指系统是否真的存储数据，而是看计算当前输出时\red{是否需要当前时刻以外的输入}。因而 $y(t)=t x(t)$ 虽然显含时间，仍然无记忆。

\begin{intuition}
  把系统看成逐点计算器：若每个输出点都能只用正下方的输入点算出，系统便无记忆；若必须向左或向右读取其他输入点，系统就有记忆。\red{无记忆性关注的是时间依赖范围}，与公式对幅值做平方、指数等非线性运算无关。
\end{intuition}

\begin{example}[逐点非线性与延迟]
  系统
  \[
    y(t)=x^2(t)+\e^{x(t)}
  \]
  只使用当前值 $x(t)$，所以无记忆。$y(t)=x(t-1)$ 读取过去值，因而有记忆。类似地，$y[n]=x^2[n]$ 无记忆，而 $y[n]=x[2n]$ 在一般下标处读取的不是 $x[n]$，所以有记忆。
\end{example}

\begin{remark}[无记忆与因果性的关系]
  无记忆系统只读取当前输入，所以\red{无记忆系统一定因果}。反过来并不成立：延迟系统 $y(t)=x(t-1)$ 只读取过去输入，是因果系统，却有记忆。
\end{remark}

\begin{pitfall}[只检查原点或个别下标]
  对 $y[n]=x[2n]$，在 $n=0$ 时确有 $y[0]=x[0]$，但系统性质要求对任意时刻成立。只要取 $n=1$，输出就需要 $x[2]$，已经足以判定系统有记忆。触发检查是：解方程“内部输入下标等于当前输出下标”是否对\red{所有时刻恒成立}。
\end{pitfall}

\section{可逆系统}
\label{sec:invertible-systems}

\begin{definition}[可逆系统]
  \label{def:invertible-system}
  若不同输入一定产生不同输出，也就是输出信号能够唯一确定输入信号，则称该系统为可逆系统（invertible system）。此时可以在原系统之后接入逆系统，使信号恢复为原输入：
  \begin{equation}
    x\longrightarrow\boxed{\text{原系统}}\longrightarrow y
    \longrightarrow\boxed{\text{逆系统}}\longrightarrow x.
    \label{eq:inverse-system}
  \end{equation}
\end{definition}

可逆性的本质是输入输出规则\red{不能丢失区分输入所需的信息}。判断时可以沿两个方向进行：若能从一般的 $y$ 唯一解出 $x$，就给出了逆系统；若能找到 $x_1\ne x_2$ 却产生同一个输出，便可立即否定可逆性。

\begin{example}[延迟与平方]
  对 $y(t)=x(t-1)$，把输出超前 $1$ 即得
  \[
    x(t)=y(t+1),
  \]
  所以延迟系统可逆。对 $y(t)=x^2(t)$，输入 $x(t)$ 与 $-x(t)$ 产生同一输出，符号信息已经丢失，因此在全体实信号上不可逆。若事先把输入限制为非负信号，则可用 $x(t)=\sqrt{y(t)}$ 恢复；可见\red{可逆性依赖允许的信号集合}。
\end{example}

\begin{example}[累加与差分]
  对离散累加系统
  \[
    y[n]=\sum_{k=-\infty}^{n}x[k],
  \]
  相邻输出相减可得
  \[
    x[n]=y[n]-y[n-1].
  \]
  因而在累加和存在的信号类上，它的逆系统是后向差分。验证逆系统时，应\red{把恢复公式代回原关系}，而不只凭运算名称猜测。
\end{example}

\begin{pitfall}[积分常数造成的不唯一]
  若系统为 $y(t)=\dd x(t)/\dd t$，直接写 $x(t)=\int y(t)\dd t$ 会漏掉积分常数。事实上 $x(t)$ 与 $x(t)+C$ 有相同导数，所以微分系统在一般可微信号上不可逆；只有额外给定一个初值，才可能唯一恢复输入。看到积分或微分时，应检查\red{边界条件是否足够}。
\end{pitfall}

\section{稳定系统}
\label{sec:stable-systems}

\begin{definition}[BIBO 稳定系统]
  \label{def:bibo-stable-system}
  若每一个有界输入都产生有界输出，则称系统为有界输入有界输出稳定系统（bounded-input bounded-output stable system），简称 BIBO 稳定系统。连续时间的判据可写成
  \[
    |x(t)|\le M_x<\infty,\ \forall t
    \quad\Longrightarrow\quad
    |y(t)|\le M_y<\infty,\ \forall t,
  \]
  其中 $M_y$ 可以依赖输入及其界 $M_x$，但不能随 $t$ 增长。离散时间情形只需把 $t$ 换成 $n$。
\end{definition}

证明稳定性时，需要从任意输入界 $M_x$ 推出一个与时间无关的有限输出界；否定稳定性时，只需找到一个有界输入使输出无界。也就是说，稳定性检验遵循\red{证明要一般、否定用反例}的逻辑。

\begin{example}[指数与多项式逐点系统]
  对实输入系统 $y(t)=\e^{x(t)}$，若 $|x(t)|\le M_x$，则
  \[
    |y(t)|=\e^{x(t)}\le \e^{|x(t)|}\le \e^{M_x},
  \]
  因而系统稳定。对
  \[
    y(t)=x^3(t)-2x^2(t)+x(t)+1,
  \]
  三角不等式给出
  \[
    |y(t)|\le M_x^3+2M_x^2+M_x+1,
  \]
  右端是有限常数，所以系统也稳定。逐点非线性并不必然破坏稳定性。
\end{example}

\begin{example}[微分、积分与差分]
  理想微分系统不是 BIBO 稳定系统：有界阶跃输入的导数包含单位冲激，不是有界的普通函数。积分器也不稳定；取有界输入 $x(t)=\stepfun(t)$，则
  \[
    y(t)=\int_{-\infty}^{t}\stepfun(\tau)\dd\tau=t\stepfun(t),
  \]
  随时间无界增长。相反，离散差分系统满足
  \[
    |x[n]-x[n-1]|\le |x[n]|+|x[n-1]|\le2M_x,
  \]
  所以稳定。这里决定稳定性的不是“积分还是微分”这类名称，而是系统能否把有界变化\red{累积或放大为无界输出}。
\end{example}

\begin{pitfall}[用一个表现良好的输入证明稳定]
  某个有界输入得到有界输出，只能说明这一个输入没有触发问题，不能证明系统稳定。正确做法是推导对所有有界输入都成立的统一上界。反之，一个有界输入产生无界输出便足以否定稳定性。
\end{pitfall}
